\section{An exact falsification of the curvature-to-real-zero inference}
\label{sec:curvatureobstruction}

The global radial-curvature theorem for the genuine Jacobi-theta
kernel provides a common \emph{bounded} Fourier-side Laguerre
positivity interval. It cannot on its own be extended to
the entire real Fourier axis. The following exact example
makes this logical boundary decisive.

\begin{theorem}[Gaussian--hyperbolic counterexample]
\label{thm:coshcounter}
Let
\[
K(u)=e^{-u^2/2}\left(1+\frac12\cosh u\right),
\qquad V(u)=-\log K(u).
\]
Then $K$ is positive, even, entire, and rapidly decreasing
on the real axis, with
\begin{align}
V''(u)
&=\frac{\cosh^2u+2\cosh u+3}
        {(\cosh u+2)^2}\ge\frac23,\\
V'''(u)
&=\frac{2(\cosh u-1)\sinh u}
        {(\cosh u+2)^3}>0\quad(u>0).
\label{eq:coshcurv}
\end{align}
Nevertheless its Fourier transform is
\begin{equation}
F(z)=\int_{\mathbb R}K(u)e^{izu}\,du
=\sqrt{2\pi}e^{-z^2/2}(1+d\cos z),
\qquad d=\frac{e^{1/2}}2,
\label{eq:coshfourier}
\end{equation}
where $1/2<d<1$. Thus $F(x)>0$ for \emph{every}
real $x$, but $F$ has nonreal zeros
\begin{equation}
z=(2k+1)\pi\pm i\operatorname{arcosh}(1/d),
\qquad k\in\mathbb Z.
\end{equation}
Its first Fourier-side Laguerre inequality is
already violated at the explicit real point $x=\pi$:
\begin{equation}\label{eq:coshlaguerre}
\boxed{
L_1[F](\pi)
=2\pi e^{-\pi^2}(1-d)(1-2d)<0.
}
\end{equation}
\end{theorem}

\begin{proof}
Put $C=\cosh u$, $S=\sinh u$. The elementary identities
$C^2-S^2=1$, $C'=S$, and $S'=C$ yield
\[
V''=1-\frac{2C+1}{(C+2)^2}
\quad\text{and}\quad
V'''=\frac{2(C-1)S}{(C+2)^3}.
\]
The strong convexity bound is exactly
\[
3(C^2+2C+3)-2(C+2)^2
=(C-1)^2\ge0.
\]
For $u>0$, $C>1$ and $S>0$, proving
strict radial curvature.

The standard Gaussian Fourier integral
$\int_{\mathbb R}e^{-u^2/2}e^{izu}du
=\sqrt{2\pi}e^{-z^2/2}$ holds for complex $z$
by analytic continuation.
The hyperbolic factor is
$\cosh u=(e^u+e^{-u})/2$; shifting the Fourier
argument by $\mp i$ and completing squares
yields \eqref{eq:coshfourier}.
Elementary exponential series show
$1<e^{1/2}<2$, hence $1/2<d<1$.
For real $x$, $1+d\cos x\ge1-d>0$.
On the other hand, at
$z=(2k+1)\pi+iy$ one has
$\cos z=-\cosh y$, and
$\cosh y=1/d>1$ has a unique $y>0$,
proving the explicit nonreal zero pair.

Finally differentiating
$F(z)=C_0e^{-z^2/2}H(z)$,
$H(z)=1+d\cos z$, gives
\[
(F')^2-FF''
=C_0^2e^{-z^2}
\bigl[H^2+(H')^2-HH''\bigr].
\]
Substituting $z=\pi$, where
$H=1-d$, $H'=0$, $H''=d$, gives
\eqref{eq:coshlaguerre}. Every sign follows
from $1/2<d<1$.
\end{proof}

\begin{proposition}[A sharp zero-location transition
inside the radial-curvature class]\label{prop:coshfamily}
For any $c\ge1/2$ define
\[
K_c(u)=e^{-u^2/2}(1+c\cosh u),\qquad
d_c=c\,e^{1/2}.
\]
All these kernels have $(-\log K_c)'''(u)>0$
for $u>0$, and their potentials are
strictly convex. Their Fourier transforms are
\[
F_c(z)=\sqrt{2\pi}e^{-z^2/2}(1+d_c\cos z).
\]
The transform has only real zeros exactly when
$d_c\ge1$, and nonreal zeros when $0<d_c<1$.
Thus the continuous curvature-positive family
passes through a real-zero transition at
$c=e^{-1/2}$, a parameter strictly greater than $1/2$.
\end{proposition}
\begin{proof}
Write $M=1+c\cosh u$ and $C=\cosh u$.
Direct differentiation yields
\[
(\log M)'''=
\frac{c\sinh u(1-cC-2c^2)}
{(1+cC)^3}.
\]
For $c\ge1/2$ and $u>0$, one has
$cC+2c^2>1$, so $(-\log K_c)'''>0$.
Moreover
\[
(\log M)''=
\frac{cC+c^2}{(1+cC)^2}<1,
\]
since the difference between denominator
and numerator is $1+cC+c^2(C^2-1)>0$.
Thus $(-\log K_c)''>0$.
Gaussian integration gives the transform.
If $d_c<1$, $\cos z=-1/d_c<-1$ has
nonreal solutions and no real solutions.
If $d_c\ge1$, solving
$\cos z=-1/d_c\in[-1,0)$ shows all
solutions are real (the boundary case
$d_c=1$ has double real zeros).
\end{proof}

\paragraph{Implications for RH.}
The counterexample proves that positivity of
$V'''$, even coupled to uniform strong log-concavity,
cannot imply either $L_1[F](x)\ge0$ for all real $x$
or the real-zero property. This does not weaken
Theorem~\ref{thm:globalcurvature} for the \emph{specific}
theta kernel; it establishes that a successful RH proof
must exploit additional theta/modular information
\emph{at the Fourier-transform level}.
A precise future target is a signed-cosine
correlation inequality
\[
\int_{\mathbb R}
G_n(w)\cos(xw)\,dw\ge0\qquad
(\forall n\ge1,\ \forall x\in\mathbb R),
\]
where
$G_n(w)=\int(2u-w)^{2n}\Phi(u)\Phi(w-u)\,du$
with the positive factorial normalization
specified earlier. The mere fact $G_n(w)\ge0$
does not provide this Fourier-positivity statement.
This distinction is consistent with the earlier
Wronskian/correlation framework of
Dimitrov--Xu~\cite{DX}.
